\chapter{Introduction}
\label{chap:introduction}

\section{Context and Motivation}

Clinical research increasingly relies on information collected during routine
patient care. In an \gls{icu}, a single admission can produce demographic
information, repeated vital-sign measurements, laboratory results, diagnoses,
and an eventual clinical outcome. Combined as \gls{ehr}, these observations
are valuable for studying patient populations and developing predictive
models. Their sensitivity, however, means that access and reuse must be
carefully controlled. This can make suitable research datasets difficult to
obtain, while uncommon outcomes may be represented by relatively few patients
even in an available dataset.

Synthetic records provide one possible complement to restricted real data. A
model is fitted to an observed cohort and then sampled to create new,
artificial records that approximate characteristics of that cohort. Research
on synthetic \gls{ehr} has demonstrated the potential of this approach for
producing realistic health data without directly distributing the source
records \cite{Yoon2023_EHRSafe}. This does not make synthetic data inherently
anonymous, clinically valid, or useful. Those properties depend on what the
model has learned and must be examined empirically.

The \gls{gan} framework learns through competition between a generator and a
discriminator \cite{Goodfellow2014GAN}. Its conditional formulation,
\glspl{cgan}, can incorporate additional information into the generation
process \cite{Mirza2014CGAN}. CTGAN and CTAB-GAN+ extend this idea to tabular
data containing categorical and continuous variables
\cite{Xu2019CTGAN,Zhao2024CTABGANPlus}. DP-CGANS was proposed for conditional
generation with differential privacy in health-data settings
\cite{Sun2023DPCGANS}.

\section{Research Problem and Objective}

Overall statistical similarity does not show whether a generator covers the
full variability of the real data. Alaa et al. therefore distinguish fidelity
from diversity, which measures the extent to which real observations are
covered by the generated distribution \cite{Alaa2022FaithfulSyntheticData}.
This distinction is especially relevant for imbalanced tabular data. A
generator can omit a minor category while causing only a small change in the
overall distribution, and the scarcity of minority observations provides
fewer opportunities to learn these cases during training
\cite{Xu2019CTGAN}. Long-tailed continuous variables present a related
difficulty because values far from the main body of the distribution can be
challenging to encode accurately \cite{Zhao2024CTABGANPlus}. In an evaluation
of severely imbalanced clinical data, Kim et al. similarly found that
conventional single-metric assessment could miss minority-class instability
and that numerical agreement did not establish equivalent clinical behaviour
\cite{Kim2026IntegratedSyntheticClinicalEvaluation}. Synthetic data should
therefore be evaluated separately for overall fidelity, rare-case coverage,
and downstream predictive utility.

Finding a discrepancy is also different from explaining it. Distributional
metrics quantify disagreement but do not necessarily identify the variables
that make synthetic records recognisable or the regions that are missing.
Methods for explaining \gls{ai} models can provide feature-level evidence.
SHapley Additive exPlanations (SHAP), for example, attribute a model output to
its input features \cite{Lundberg2017SHAP}. Explanations have previously been
used to analyse or guide generative models
\cite{Nagisetty2020XAIGAN,Rozendo2024XGAN}, and sample reweighting has been
explored as a way to improve generated distributions
\cite{Diefenbacher2020DCTRGAN}. For tabular health data, however, there remains
a practical gap between identifying an important feature and deciding which
real training records should be emphasised.

The objective of this thesis is to examine whether that gap can be bridged by
combining feature explanations with measured underrepresentation. The
empirical study uses MIMIC-IV \cite{PhysioNet-mimiciv-3.1} and WiDS, two
tabular intensive-care datasets, and compares CTAB-GAN+, CTGAN, and DP-CGANS.
Hospital mortality provides the main rare-outcome prediction task.

DP-CGANS is run with private training disabled in this study. It consequently
serves as an architectural comparison model rather than as evidence of a
formal differential-privacy guarantee. The privacy-related evaluation is
limited to indicators of memorisation, such as exact matches and
nearest-neighbour distances.

\section{Research Questions}

The study is organised around the following questions:

\begin{enumerate}
    \item \textbf{RQ1:} How well do CTAB-GAN+, CTGAN, and DP-CGANS represent
    global distributions, tails, rare cases, and downstream predictive
    utility in tabular intensive-care data?

    \item \textbf{RQ2:} What do post-hoc TreeSHAP explanations of a
    real-versus-synthetic detector and GradientSHAP explanations of the
    internal CTAB-GAN+ discriminator reveal about detectable synthetic-data
    artifacts, and to what extent do their feature rankings agree?

    \item \textbf{RQ3:} Can explanation-guided reweighting improve the
    fidelity, rare-case representation, or downstream utility of synthetic
    data compared with an unmodified generator and a uniform-augmentation
    control, and what trade-offs arise?
\end{enumerate}

\section{Study Overview}

The first part of the study establishes a common baseline for the three
generator architectures. Each model is trained and assessed under the same
data-partitioning and reporting procedure. The comparison considers overall
statistical fidelity, representation of tails and rare cases, transfer to a
real mortality-prediction task, real--synthetic distinguishability, and
memorisation-risk indicators.

The second part tests a targeted retraining intervention. A classifier is
trained after baseline generation to distinguish held-out real records from
synthetic records. TreeSHAP is used to determine which features contribute to
correctly identifying synthetic examples. These feature priorities are
combined with distribution and tail deficits that indicate where synthetic
probability mass is lower than real probability mass. Real training rows in
those regions are sampled more frequently when constructing an augmentation
set. An unmodified baseline and an equally sized uniform augmentation provide
controls for judging whether any change is attributable to the guidance
signals rather than simply to adding rows.

The third part examines the CTAB-GAN+ discriminator itself. GradientSHAP is
applied to discriminator snapshots collected during training, allowing its
feature attributions to be followed over time and compared with those of the
post-hoc detector. This is an independent explanation experiment: its
attributions are not used to calculate sampling weights or retrain any
generator.

Together, these experiments provide a rare-case-focused comparison of tabular
generators and a controlled test of explanation-guided data reweighting. The
central methodological contribution is the separation of two roles:
feature-level explanations indicate where an audit should focus, while
region-level deficits determine which real records receive additional
sampling probability.

\section{Thesis Outline}

The next chapter introduces the relevant work on tabular data generation,
synthetic health-data evaluation, and model explanation.
Chapter~\ref{chap:methodology} then presents the experimental design, the
reweighting method, and the separate internal-discriminator analysis. The
evaluation chapter defines the measures used for statistical representation,
rare cases, predictive utility, detectability, and memorisation risk. The
results and discussion chapters report and interpret the findings. The final
chapter answers the research questions and summarises the limitations and
opportunities for further work.
